% I. E. Irodov, Problems in General Physics, Part II §2.7, 246--257.
% 对照用户上传PDF页99--100，印刷页90--91；简答PDF页304 (printed p295).
% \kappa 为热导率，\eta 为动力黏度，热流量密度 j_Q 取向外正。
\thermotransportproblem{246}{90}
利用泊肃叶方程，求在一根长为 $\ell$、内半径为 $a$
的细管中，管两端保持气体压强 $p_1,p_2$ 时，
每单位时间通过管道的气体质量 $\dot m$。
气体为摩尔质量 $M$ 的理想气体，温度恒为 $T$、
动力黏度为 $\eta$。
\exsolution
气体可压缩，沿管道压强变化，因此直接把
不可压缩泊肃叶体积流量 $Q$ 乘入口密度并不正确。
在截面局部采用泊肃叶定律
$Q(x)=-(\pi a^4/8\eta)\mathrm dp/\mathrm dx$，
质量流量为
\[
\dot m=\rho(x)Q(x)
=-\frac{\pi a^4M}{8\eta RT}\,
p(x)\frac{\mathrm dp}{\mathrm dx},
\]
其中用到了等温理想气体关系
$\rho=pM/(RT)$。稳态时 $\dot m$ 与位置无关，
积分从 $x=0$ 到 $\ell$ 得
\[
\dot m\ell
=-\frac{\pi a^4M}{8\eta RT}
\int_{p_1}^{p_2}p\,\mathrm dp
=\frac{\pi a^4M}{16\eta RT}(p_1^2-p_2^2).
\]
因此质量流率的大小为
\[
\boxed{\dot m=\frac{\pi a^4 M}{16\eta RT\ell}
|p_1^2-p_2^2|}.
\]
\excheck 极限 $p_1\to p_2$ 时流率趋零；
当压强差远小于平均压强时，
才可将上式近似为平均密度乘普通泊肃叶流量。

\thermotransportproblem{247}{90}
一根侧面完全绝热的复合导热棒，
由长度 $\ell_1,\ell_2$、热导率
$\kappa_1,\kappa_2$ 的两段串接。
两端分别恒温 $T_1,T_2$。
求两段交界处的稳态温度。
\exsolution
两段截面积相同，稳态下热流密度相等：
\[
\kappa_1\frac{T_1-T_*}{\ell_1}
=\kappa_2\frac{T_*-T_2}{\ell_2}.
\]
两边移项，得到
$(\kappa_1/\ell_1+\kappa_2/\ell_2)T_*
=\kappa_1T_1/\ell_1+\kappa_2T_2/\ell_2$，
故
\[
\boxed{T_*=
\frac{\kappa_1T_1/\ell_1+\kappa_2T_2/\ell_2}
{\kappa_1/\ell_1+\kappa_2/\ell_2}}.
\]
\excheck 当 $\kappa_1/\ell_1\gg\kappa_2/\ell_2$ 时
$T_*\to T_1$，符合第一段热阻很小的极限。

\thermotransportproblem{248}{90}
两根截面积相同、长度为 $\ell_1,\ell_2$、
热导率为 $\kappa_1,\kappa_2$ 的棒首尾相接，
侧面绝热。求与之具有相同稳态导热能力、
总长为 $\ell_1+\ell_2$ 的等效均匀棒的热导率。
\exsolution
等截面串联的热阻相加：
\[
\mathcal R_{\rm th}=
\frac{\ell_1}{\kappa_1 S}+\frac{\ell_2}{\kappa_2 S}
=\frac{\ell_1+\ell_2}{\kappa_{\rm eq}S}.
\]
解得
\[
\boxed{\kappa_{\rm eq}=
\frac{\ell_1+\ell_2}
{\ell_1/\kappa_1+\ell_2/\kappa_2}}.
\]
\excheck 是按长度加权的\emph{调和平均}，
而非算术平均；原版书后公式的 PDF 文本抽取
可能遗漏分式横线，应以热阻相加为准。

\thermotransportproblem{249}{91}
长为 $\ell$ 的均匀导热棒侧面绝热，
材料热导率与温度的关系为
$\kappa(T)=\alpha/T$，
其中 $\alpha>0$。
坐标 $x=0,\ell$ 两端分别维持
$T_1,T_2>0$。求 $T(x)$ 与稳态热流密度。
\exsolution
一维稳态傅里叶定律
$j_Q=-\kappa(T)\mathrm dT/\mathrm dx$
给出
\[
j_Q=-\frac{\alpha}{T}\frac{\mathrm dT}{\mathrm dx}
=-\alpha\frac{\mathrm d(\ln T)}{\mathrm dx}.
\]
$j_Q$ 恒定，故 $\ln T$ 随 $x$ 线性变化：
\[
\ln T(x)=\ln T_1+
\frac{x}{\ell}\ln\frac{T_2}{T_1}.
\]
因此
\[
\boxed{T(x)=T_1\left(\frac{T_2}{T_1}\right)^{x/\ell}},
\qquad
\boxed{j_Q=\frac{\alpha}{\ell}\ln\frac{T_1}{T_2}}.
\]
\excheck 若 $T_1>T_2$，热量沿 $+x$ 方向输运，
$j_Q>0$；原书简答以相反方向定义正热流时
会写成 $\ln(T_2/T_1)$，物理结果相同。

\thermotransportproblem{250}{91}
两块金属的总热容分别为 $C_1,C_2$，
以一根长 $\ell$、截面积 $S$、
热导率 $\kappa$ 很小的棒连接；
整个系统与外界绝热，忽略导热棒自身热容。
初始温差为 $\Delta T_0$。
求金属块温差随时间的变化。
\exsolution
记两块瞬时温度 $T_1(t),T_2(t)$，
热功率 $P=(\kappa S/\ell)(T_1-T_2)$。
能量守恒与各块热平衡方程为
\[
C_1\dot T_1=-P,\qquad C_2\dot T_2=P.
\]
两式相减：
\[
\frac{\mathrm d}{\mathrm dt}(T_1-T_2)
=-\frac{\kappa S}{\ell}
\left(\frac1{C_1}+\frac1{C_2}\right)(T_1-T_2).
\]
积分得
\[
\boxed{\Delta T(t)=\Delta T_0e^{-t/\tau}},\quad
\boxed{\tau^{-1}=\frac{\kappa S}{\ell}
\left(\frac1{C_1}+\frac1{C_2}\right)}.
\]
\excheck 两块总体能量恒定，最终温度相同；
$\tau$ 的量纲为秒。

\thermotransportproblem{251}{91}
两块大平行板间距 $\ell$，
分别恒温 $T_1,T_2$；
板间均匀材料的热导率满足
$\kappa(T)\propto\sqrt{T}$。
求稳态温度分布 $T(x)$，
以 $T_1$ 侧为 $x=0$。
\exsolution
\textbf{按原版图像恢复根号：}
题面是 $\kappa\propto\sqrt{T}$，
纯文本抽取会丢失 $\sqrt{\ }$。
写成 $\kappa(T)=A\sqrt T$，
傅里叶定律给
$j_Q=-A\sqrt T\,T'(x)$，
故 $T^{3/2}$ 随 $x$ 线性变化：
\[
\frac{\mathrm d(T^{3/2})}{\mathrm dx}
=-\frac{3j_Q}{2A}=\text{常数}.
\]
按端点条件，
\[
T(x)^{3/2}
=T_1^{3/2}+\frac{x}{\ell}
(T_2^{3/2}-T_1^{3/2}).
\]
故
\[
\boxed{T(x)=T_1
\left[1+\frac{x}{\ell}
\left(\left(\frac{T_2}{T_1}\right)^{3/2}-1\right)
\right]^{2/3}}.
\]
\excheck 端点 $x=0,\ell$ 分别得到 $T_1,T_2$；
若漏掉根号将错误得到平方根型温度解。

\thermotransportproblem{252}{91}
两块大水平板之间为氦气，间距
$\ell=50\,\mathrm{mm}$，
下板温度 $T_1=290\,\mathrm K$，
上板温度 $T_2=330\,\mathrm K$，
气压接近标准大气压。
求稳态热流密度。
\exsolution
在连续气体传热模型中，
单原子气体的热导率
$\kappa\simeq\frac13\rho c_V\bar v\lambda$，
其中 $\bar v=\sqrt{8RT/(\pi M)}$，
$\lambda=(\sqrt2\pi d^2n)^{-1}$、
$\rho c_V=\frac32nk_B$。
代入整理：
\[
\kappa(T)\simeq
\frac{k_B}{\pi d^2}\sqrt{\frac{RT}{\pi M}}
=A\sqrt T.
\]
由傅里叶定律积分，
朝向低温下板的热流密度大小为
\[
|j_Q|=\frac1\ell\int_{T_1}^{T_2}\kappa(T)\,\mathrm dT
=\boxed{\frac{2k_B}{3\pi^{3/2}d^2\ell}
\sqrt{\frac RM}(T_2^{3/2}-T_1^{3/2})}.
\]
按氦气有效碰撞直径
$d\simeq0.20\,\mathrm{nm}$、
$M\simeq0.004\,\mathrm{kg/mol}$，
\[
\boxed{|j_Q|\simeq40\,\mathrm{W/m^2}}.
\]
\excheck 原题未给 $d$，数值依赖采用的
氦气有效硬球碰撞直径；
热流从330 K上板流向290 K下板。

\thermotransportproblem{253}{91}
两大平板间距 $\ell=5.0\,\mathrm{mm}$，
板间为压强 $p=1.0\,\mathrm{Pa}$ 的氦气。
两板温度分别为 $17^\circ\mathrm C$、
$37^\circ\mathrm C$。
求氦原子平均自由程与热流密度。
\exsolution
平均温度约
$T=(290+310)/2\simeq300\,\mathrm K$。
硬球自由程
\[
\lambda=\frac{k_BT}{\sqrt2\pi d^2p}.
\]
取氦有效直径 $d\simeq0.20\,\mathrm{nm}$，
得
\[
\boxed{\lambda\simeq23\,\mathrm{mm}>\ell}.
\]
故气体处于自由分子输运而非连续傅里叶传热区。
采用原书的自由分子通量近似
\[
|j_Q|=\frac{p\bar v(T_2-T_1)}
{6T(\gamma-1)},\qquad
\bar v=\sqrt{\frac{8RT}{\pi M}},
\]
单原子氦气 $\gamma=5/3$，
$M=0.004\,\mathrm{kg/mol}$，得
\[
\boxed{|j_Q|\simeq22\,\mathrm{W/m^2}}.
\]
\excheck 比较 $\lambda$ 与板间距
决定使用哪一种模型；
若将连续介质傅里叶公式外推到
$\lambda>\ell$，不符合本题设定。

\thermotransportproblem{254}{91}
内外半径为 $R_1<R_2$ 的两根同轴无限长
圆筒之间充满均匀导热物质，
其热导率为常数。两筒分别恒温
$T_1,T_2$。求稳态温度分布 $T(r)$。
\exsolution
取长度 $L$ 的柱形等温面，
稳态径向热功率不随半径变化：
\[
P=-2\pi rL\kappa\frac{\mathrm dT}{\mathrm dr}.
\]
因此 $rT'(r)=\text{常数}$，
积分有 $T(r)=A\ln r+B$。
代入 $T(R_1)=T_1$、
$T(R_2)=T_2$ 得
\[
\boxed{T(r)=T_1+
\frac{T_2-T_1}{\ln(R_2/R_1)}
\ln\frac r{R_1}},\qquad R_1\le r\le R_2.
\]
\excheck $T(r)$ 对 $\ln r$ 线性，
不是对半径 $r$ 线性。

\thermotransportproblem{255}{91}
把前题的两同轴圆筒改为
内外半径 $R_1<R_2$ 的同心球壳，
球面分别恒温 $T_1,T_2$。
求稳态温度分布。
\exsolution
半径 $r$ 的球面上的传热功率为
\[
P=-4\pi r^2\kappa\,T'(r)=\text{常数}.
\]
故 $r^2T'(r)$ 恒定，
积分得到 $T(r)=A+B/r$。
代入两边界值，
\[
\boxed{T(r)=T_1+
(T_2-T_1)
\frac{1/R_1-1/r}{1/R_1-1/R_2}},
\qquad R_1\le r\le R_2.
\]
\excheck 分子在 $r=R_1$ 为零，
在 $r=R_2$ 等于分母，端点温度无误。

\thermotransportproblem{256}{91}
截面半径为 $R$ 的均匀长导线，
内部有恒定电流，使每单位体积产生
恒定焦耳热功率 $w$。
导线热导率为 $\kappa$，
其外表面维持温度 $T_0$。
求截面内的稳态温度分布。
\exsolution
轴对称径向稳态导热方程为
\[
\frac1r\frac{\mathrm d}{\mathrm dr}
\left(r\frac{\mathrm dT}{\mathrm dr}\right)
=-\frac w\kappa.
\]
积分一次：
$rT'(r)=-wr^2/(2\kappa)+C$。
由于轴线上温度梯度有限，$C=0$。
再积分：
$T(r)=-wr^2/(4\kappa)+C_2$。
由 $T(R)=T_0$ 得
\[
\boxed{T(r)=T_0+\frac w{4\kappa}(R^2-r^2)}.
\]
\excheck 中心最热，$T(0)-T_0=wR^2/(4\kappa)$；
正则性条件排除了轴线附近的 $\ln r$ 项。

\thermotransportproblem{257}{91}
半径为 $R$ 的均匀实心球，
内部每单位体积均匀产生热功率 $w$。
热导率为 $\kappa$、球表面温度为 $T_0$。
求球内稳态温度分布。
\exsolution
球对称稳态传热方程为
\[
\frac1{r^2}\frac{\mathrm d}{\mathrm dr}
\left(r^2\frac{\mathrm dT}{\mathrm dr}\right)
=-\frac w\kappa.
\]
积分一次
$r^2T'(r)=-wr^3/(3\kappa)+C$；
由中心有限温度梯度得 $C=0$。
积分并施加表面边界 $T(R)=T_0$，
得到
\[
\boxed{T(r)=T_0+\frac w{6\kappa}(R^2-r^2)}.
\]
\excheck 长圆柱（第256题）系数为 $1/4$，
球体为 $1/6$，反映不同径向通量面积。
